% part: intuitionistic-logic
% chapter: propositions-as-types
% section: type-preservation

\documentclass[../../../include/open-logic-section]{subfiles}

\begin{document}

\olfileid{int}{pty}{tpr}

\olsection{类型保持}

前文所考察的归约与排列转换，同样可以直接在系统 \Log{tN2i} 和 \Log{tN3i} 中定义。在系统 \Log{tN3i} 中，$\lambd$-项~$N$ 相对于语境~$\Gamma$ 的\emph{类型}被定义为满足 $\Gamma \Sequent N : !A$ 在 \Log{tN3i} 中有推导的公式~$!A$。$\lambd$-项的 $\beta$-归约规则与推导的归约及排列转换精确对应，这一事实有一个重要推论：$\lambd$-项相对于语境~$\Gamma$ 的类型在 $\beta$-转换后保持不变。这一性质称为\emph{类型保持}。

\begin{prop}
若 $N$ 相对于 $\Gamma$ 具有类型~$!A$ 并且 $N \redone N'$，则 $N'$ 相对于~$\Gamma$ 也具有类型~$!A$。
\end{prop}

\begin{proof}
通过对~$N$ 以及 $\Gamma \Sequent N : !A$ 的推导~$\delta$ 的高度施归纳，我们容易确立以下事实：若 $M$ 是~$N$ 的子项，则 $\delta$ 包含一个结束于 $\Gamma' \Sequent M : !B$ 的子推导~$\delta_1$。若 $M$ 是可约式，则可对 $\delta_1$ 施加某种转换。对该 $\delta_1$ 施加归约，我们便得到 $\Gamma' \Sequent M' : !B$ 的推导~$\delta_1'$。审查 \Log{tN3i} 的转换规则可知，$M'$ 即为~$M$ 的归约结果。将 $\delta$ 中的 $\delta_1$ 替换为 $\delta_1'$，便得到 $\Gamma \Sequent N' : !A$ 的推导~$\delta'$，其中 $N'$ 是将 $N$ 中的子项~$M$ 替换为~$M'$ 所得。

例如，考虑 \Intro{\land} 后接 \Elim{\land} 的归约转换：
\[
  \bottomAlignProof
  \AxiomC{}
  \Deduce$\Gamma \fCenter N_1 : !B_1$
  \AxiomC{}
  \Deduce$\Gamma \fCenter N_2 : !B_2$
  \RightLabel{\Intro{\land}}
  \BinaryInf$\Gamma \fCenter \pair{N_1}{N_2} : !B_1 \land !B_2$
  \RightLabel{\Elim{\land}[i]}
  \UnaryInf$\Gamma \fCenter \proj{i}{\pair{N_1}{N_2}} : !B_i$
  \DisplayProof
  \quad\redone\quad
  \bottomAlignProof
  \AxiomC{}
  \Deduce$\Gamma \fCenter N_i : !B_i$
  \DisplayProof
  \]
我们看到右侧推导（即 $\delta_1'$）的证明项是 $N_i$。这恰是对左侧推导~($\delta_1$) 的证明项 $\proj{i}{\pair{N_1}{N_2}}$ 施行 $\beta$-转换的结果。

又如，考虑一个 \Intro{\lif} 后接 \Elim{\lif} 推理，以及对其施加的归约转换：
\begin{multline*}
  \bottomAlignProof
  \AxiomC{$\delta_2$}
  \Deduce$x: !B_1, \Gamma \fCenter N : !B_2$
  \RightLabel{\Intro{\lif}}
  \UnaryInf$\Gamma \fCenter \lambd[\typeof{x}{!B_1}][N] : !B_1 \lif !B_2$
  \AxiomC{}
  \RightLabel{$\delta_3$}
  \Deduce$\Gamma \fCenter M : !B_1$
  \RightLabel{\Elim{\lif}}
  \BinaryInf$\Gamma \fCenter (\lambd[\typeof{x}{!B_1}][N] M) : !B_2$
  \DisplayProof
  \quad\redone\\
  \bottomAlignProof
  \AxiomC{}
  \RightLabel{$\delta_3$}
  \Deduce$\Gamma \fCenter M : !B_1$
  \RightLabel{$\Subst{\delta_2}{\delta_3}{x:!B_1}$}
  \Deduce$\Gamma \fCenter \Subst{N}{M}{x} : !B_2$
  \DisplayProof
  \end{multline*}
右侧推导的证明项同样恰是对左侧推导的证明项进行 $\beta$-归约的结果。当然，我们必须仔细验证（通过对~$\delta_2$ 的高度施归纳）：若在~$\delta_2$ 中将所有形如 $\Gamma' \Sequent x:!B_1$ 的公理替换为 $\Gamma \Sequent M:!B_1$ 的推导~$\delta_3$，确实会得到 $\Gamma \Sequent
  \Subst{N}{M}{x} : !B_2$ 的推导 $\Subst{\delta_2}{\delta_3}{x:!B_1}$。
\end{proof}

推导与 $\lambd$-项之间的紧密对应还有另一个推论。若 $N$ 是相对于语境~$\Gamma$ 具有类型 $!A$ 的 $\lambd$-项，且包含一个可约式~$M$（即它不是范式），则与之对应的 \Log{tN3i} 证明（即 $\Gamma \Sequent N : !A$ 的推导~$\delta$）包含一个可施加转换的子推导（该子推导结束于 $\Gamma' \Sequent M: !B$）。若通过将 $M$ 替换为其归约结果使得 $N \redone N'$，则对~$\delta$ 施加相应归约的结果便是 $\Gamma \Sequent N' : !A$ 的推导。因此，每一个不在范式中的 $\lambd$-项都能 $\beta$-归约到另一个项。这一性质称为\emph{进展性}。

进展性与类型保持共同保证了 $\lambd$-项的 $\beta$-归约永远不会“卡住”：若一个项是良类型的且不在范式中，它就会归约到另一个良类型的项（且该项与原项类型相同）。

\end{document}